% TOP_PROBLEM: {"id": 4900014, "title": "Quantum unique ergodicity", "category_id": 11, "category": {"id": 11, "name": "dynamical_systems", "display_name": "Dynamical Systems", "description": "Problems about long-term behavior of deterministic systems, Hamiltonian dynamics, and periodic orbits.", "slug": "dynamical-systems", "order_index": 11, "created_at": "2026-07-31T15:26:25.671Z"}, "status": "partially_solved"}
% FIRST_POSTED: null
% ENTRY_KIND: statement_only
% Imported from: batch16.tex; UnsolvedMath ID 4900014
\documentclass[11pt]{article}
\usepackage[margin=25mm]{geometry}
\usepackage{fontspec}
\setmainfont{Latin Modern Roman}
\usepackage{amsmath,amssymb,mathtools}
\usepackage{unicode-math}
\setmathfont{Latin Modern Math}
\usepackage{xurl,hyperref}
\hypersetup{colorlinks=true,urlcolor=blue}
\setcounter{secnumdepth}{0}
\setlength{\parindent}{0pt}
\setlength{\parskip}{5pt}
\setlength{\emergencystretch}{3em}
\newcommand{\sref}[2]{\hyperlink{source:#1:#2}{[S#2]}}
\newcommand{\cataloguescope}[1]{\paragraph{Recorded target.}#1}
\begin{document}
% UnsolvedMath ID 4900014; source catalogue code AMR-048-0014
\setcounter{section}{4900013}
\section{Quantum unique ergodicity}\label{4900014}
{\small\sffamily UnsolvedMath ID 4900014 · Dynamical Systems}
\subsection{Definitions and mathematical statement}

\paragraph{Shared notation.}$\mathbb N=\{1,2,\ldots\}$, and $\mathbb Z,\mathbb Q,\mathbb R,\mathbb C$ denote the integers, rationals, reals and complex numbers. Any other conventions are specified locally.


Let $(M,g)$ be a closed connected smooth Riemannian manifold of dimension at least two, with strictly negative sectional curvature. Write $d\operatorname{vol}_g$ for its volume measure, $\operatorname{Vol}_g(M)=\int_M d\operatorname{vol}_g$, and $-\Delta_g$ for the nonnegative Laplace--Beltrami operator. An orthonormal sequence of eigenfunctions $(\varphi_j)_{j\geq1}$ satisfies
\[
-\Delta_g\varphi_j=\lambda_j\varphi_j,\qquad
\int_M\varphi_j\overline{\varphi_k}\,d\operatorname{vol}_g=\delta_{jk},\qquad \lambda_j\longrightarrow\infty,
\]
where $\delta_{jk}$ is one for $j=k$ and zero otherwise. Define the probability measure $\mu_j=|\varphi_j|^2d\operatorname{vol}_g$. Weak convergence to normalized volume means convergence against every continuous test function $h:M\to\mathbb R$.
\paragraph{Statement recorded in the supplied source.}
Does every such sequence, on every such $(M,g)$, satisfy
\[
\lim_{j\to\infty}\int_M h(x)|\varphi_j(x)|^2\,d\operatorname{vol}_g(x)
=\frac{1}{\operatorname{Vol}_g(M)}\int_M h(x)\,d\operatorname{vol}_g(x)
\quad\text{for every }h\in C(M,\mathbb R)?
\]
The question concerns all high-energy sequences, including arbitrary choices within eigenspaces. The source also formulates a stronger phase-space equidistribution conjecture; the displayed density statement is the registered target. \sref{4900014}{2}
\subsection{Short English statement}
On every compact negatively curved manifold without boundary, must the probability density of every sequence of normalized high-energy Laplace eigenfunctions become uniformly distributed over the manifold?
\subsection{Sources}
\begin{description}
\item[\hypertarget{source:4900014:1}{S1}] ULAM.AI and the UnsolvedMath Contributors, UnsolvedMath: A Curated Collection of Open Mathematics Problems, record 4900014 (AMR-048-0014). CC BY 4.0. \url{https://huggingface.co/datasets/ulamai/UnsolvedMath}.
\item[\hypertarget{source:4900014:2}{S2}] Nalini Anantharaman and Stéphane Nonnenmacher, Half-delocalization of Eigenfunctions for the Laplacian on an Anosov Manifold (2007), introduction, p. 1, and section 1, pp. 3--5. \url{https://arxiv.org/abs/math-ph/0610019}.
\end{description}
\label{end:4900014}
\end{document}
